\chapter{Alcohol Poisoning}

\section*{Definition and Overview}
Alcohol poisoning (acute alcohol intoxication) occurs when a person drinks enough alcohol in a short time to overwhelm the body's ability to break it down. Alcohol then builds up in the blood and suppresses the brain centres that control breathing, heart rate, and the gag reflex. It is a medical emergency that can cause permanent brain damage or death if not treated promptly.

Alcohol poisoning is different from a hangover or ordinary drunkenness. A person with alcohol poisoning may be unconscious, unrousable, or unable to protect their airway, and they require immediate emergency medical care.

\section*{Epidemiology}
Alcohol poisoning is one of the most common reasons for emergency department visits and can be fatal. It is most common in young adults, particularly in the setting of binge drinking, but it can affect anyone -- including experienced drinkers who drink faster than usual and people who mix alcohol with other depressant drugs.

Drinking games, shots, and very high-strength drinks increase the risk substantially. Deaths from alcohol poisoning occur every year, most often from choking on vomit or from breathing stopping while the person is asleep or unconscious.

\section*{Aetiology and Pathophysiology}
Alcohol poisoning happens when the blood alcohol concentration rises to a level that suppresses vital brain functions.

\begin{itemize}
    \item \textbf{Drinking fast:} a large amount of alcohol in a short time, faster than the liver can process it
    \item \textbf{Binge drinking:} typically defined as four or more standard drinks on one occasion
    \item \textbf{High-strength drinks:} spirits, pure alcohol, and drinks of unknown strength
    \item \textbf{Mixing with other drugs:} benzodiazepines, opioids, or sleeping tablets greatly increase the effect
    \item \textbf{Body size and tolerance:} smaller people and those new to drinking reach dangerous levels more quickly
\end{itemize}

Alcohol is a depressant. At high levels it suppresses the parts of the brain that keep the person breathing and that trigger the cough and gag reflex. It also causes blood sugar to fall, body temperature to drop, and vomiting to occur, all of which add to the danger.

\section*{Clinical Features}
\begin{itemize}
    \item Confusion, stupor, or an inability to be woken
    \item Vomiting, especially while drowsy or asleep
    \item Slow breathing (fewer than eight breaths per minute) or irregular breathing
    \item Cold, clammy, pale, or bluish skin
    \item Low body temperature (hypothermia)
    \item Slow heart rate
    \item Seizures
    \item Choking on vomit, which can block the airway or cause suffocation
\end{itemize}

\section*{Differential Diagnosis}
\begin{itemize}
    \item Head injury, such as from a fall while drinking
    \item Hypoglycaemia (low blood sugar), especially in people with diabetes
    \item Overdose of other drugs, such as benzodiazepines, opioids, or GHB
    \item Stroke or a seizure disorder
    \item Meningitis or severe infection, which can cause confusion and drowsiness
    \item Hypothermia from being exposed to the cold
\end{itemize}

\section*{When to Seek Medical Care}
Alcohol poisoning is a medical emergency. Never assume that an unconscious person will simply "sleep it off", and never leave them alone.

\textbf{Call 000 (or local emergency services) for:}
\begin{itemize}
    \item A person who is unconscious or cannot be woken
    \item Breathing that is slow, irregular, or stops
    \item Vomiting while drowsy, or any difficulty breathing
    \item Cold, clammy, pale, or bluish skin
    \item A seizure, or signs of a head injury
    \item Any doubt about whether the person is safe
\end{itemize}

While waiting for help, turn the person onto their side and keep their airway clear if they are vomiting.

\section*{Diagnosis and Investigations}
\begin{itemize}
    \item \textbf{Clinical assessment:} assessment of the level of consciousness, breathing, and vital signs
    \item \textbf{Blood alcohol test:} to confirm and measure the level of intoxication
    \item \textbf{Blood glucose test:} to check for low blood sugar
    \item \textbf{Blood tests:} for electrolytes, liver function, and other effects of alcohol
    \item \textbf{ECG:} to check the heart rhythm
    \item \textbf{Head CT:} if there is any chance of head injury
    \item \textbf{Tests for other drugs:} if a mixed overdose is suspected
\end{itemize}

\section*{Management}
Treatment in hospital is supportive care while the body clears the alcohol.

\begin{itemize}
    \item \textbf{Airway and breathing support:} oxygen, suction of vomit, and sometimes a breathing tube or ventilator
    \item \textbf{Intravenous fluids:} to treat dehydration and low blood pressure
    \item \textbf{Blood sugar correction:} glucose for low blood sugar
    \item \textbf{Warming:} treatment of hypothermia
    \item \textbf{Thiamine:} given to prevent complications in people with poor nutrition
    \item \textbf{Treatment of seizures and other complications:} as they arise
    \item \textbf{Observation:} monitoring until alcohol levels fall and the person is safe, with discharge only once fully alert
\end{itemize}

\section*{Complications}
\begin{itemize}
    \item Choking on vomit, which can block the airway
    \item Aspiration pneumonia, from breathing vomit into the lungs
    \item Respiratory arrest -- breathing stops
    \item Hypoglycaemia and seizures
    \item Hypothermia and heart rhythm problems
    \item Brain damage from lack of oxygen
    \item Death, in severe cases
\end{itemize}

\section*{Prognosis and Outcomes}
Most people who receive prompt medical care for alcohol poisoning recover fully without lasting harm. The outcome depends on how much was drunk, how quickly care is given, and whether other drugs were involved. Delays in seeking help are the main reason alcohol poisoning becomes fatal. People who survive are often given advice about safer drinking, and a hospital stay may provide an opportunity to address problem drinking.

\section*{Prevention and Self-Care}
\begin{itemize}
    \item \textbf{Pace drinking:} no more than one standard drink per hour, and eat before and during drinking
    \item \textbf{Avoid high-risk drinking:} avoid drinking games, shots, and high-strength drinks
    \item \textbf{Alternate drinks:} alternate alcoholic drinks with water or soft drinks
    \item \textbf{No mixing:} do not mix alcohol with medicines or other drugs
    \item \textbf{Never on an empty stomach:} do not drink when unwell or exhausted
    \item \textbf{Look after friends:} if anyone becomes drowsy or unresponsive after drinking, get help immediately rather than letting them sleep it off
    \item \textbf{Young people:} keep alcohol away from children and young people
\end{itemize}

\section*{Sources and Further Reading}
The following sources provide reliable, up-to-date information on alcohol poisoning.

\begin{itemize}
    \item NHS. \textit{Conditions: alcohol poisoning.} https://www.nhs.uk/conditions/
    \item MedlinePlus. \textit{Alcohol intoxication or poisoning: first aid.} https://medlineplus.gov
    \item Mayo Clinic. \textit{Alcohol poisoning: symptoms, treatment, and first aid.} https://www.mayoclinic.org
    \item Centers for Disease Control and Prevention. \textit{Alcohol poisoning and binge drinking.} https://www.cdc.gov
    \item World Health Organization. \textit{Alcohol and emergency care resources.} https://www.who.int
\end{itemize}
